\documentclass[11pt,a4paper]{article}

\usepackage[T1]{fontenc}
\usepackage[utf8]{inputenc}
\usepackage[italian]{babel}
\usepackage{lmodern}
\usepackage{microtype}
\usepackage[a4paper,margin=2.35cm,headheight=15pt]{geometry}
\usepackage{amsmath,amssymb,mathtools}
\usepackage{booktabs}
\usepackage{array}
\usepackage{enumitem}
\usepackage{graphicx}
\usepackage{tikz}
\usetikzlibrary{arrows.meta,positioning,calc,fit,backgrounds}
\usepackage{xcolor}
\usepackage{hyperref}
\usepackage{xurl}
\usepackage{fancyhdr}
\usepackage{longtable}
\usepackage{tabularx}
\usepackage{titlesec}
\usepackage{framed}
\usepackage{ragged2e}
\usepackage{listings}

\definecolor{ddtadark}{RGB}{28,38,52}
\definecolor{ddtagray}{RGB}{242,244,247}
\definecolor{ddtablue}{RGB}{42,88,135}
\definecolor{ddtagreen}{RGB}{48,111,78}
\definecolor{ddtaorange}{RGB}{148,91,25}
\definecolor{ddtared}{RGB}{140,48,48}

\hypersetup{
  colorlinks=true,
  linkcolor=ddtablue,
  urlcolor=ddtablue,
  citecolor=ddtablue,
  pdftitle={DDTA SecurityRequirement S2 - semantic closure and facial-access camera corpus},
  pdfauthor={Documentation-Driven Threat Analysis research workspace}
}

\pagestyle{fancy}
\fancyhf{}
\lhead{DDTA -- SecurityRequirement}
\rhead{S2 semantic closure}
\cfoot{\thepage}

\titleformat{\section}{\Large\bfseries\color{ddtadark}}{\thesection}{0.7em}{}
\titleformat{\subsection}{\large\bfseries\color{ddtadark}}{\thesubsection}{0.7em}{}
\titleformat{\subsubsection}{\normalsize\bfseries\color{ddtadark}}{\thesubsubsection}{0.7em}{}

\setlist[itemize]{leftmargin=1.55em,itemsep=0.25em,topsep=0.35em}
\setlist[enumerate]{leftmargin=1.75em,itemsep=0.25em,topsep=0.35em}
\setlength{\parskip}{0.35em}
\setlength{\parindent}{0pt}
\setlength{\emergencystretch}{2em}

\lstset{
  basicstyle=\ttfamily\small,
  columns=fullflexible,
  keepspaces=true,
  frame=single,
  framerule=0.3pt,
  rulecolor=\color{ddtagray},
  backgroundcolor=\color{ddtagray},
  xleftmargin=0.8em,
  xrightmargin=0.8em,
  aboveskip=0.6em,
  belowskip=0.6em,
  breaklines=true
}

\newcommand{\closed}{\textbf{CLOSED}}
\newcommand{\strongcandidate}{\textbf{STRONG CANDIDATE}}
\newcommand{\candidate}{\textbf{CANDIDATE}}
\newcommand{\openstatus}{\textbf{OPEN}}
\newcommand{\deferred}{\textbf{DEFERRED}}
\newcommand{\rejected}{\textbf{REJECTED}}
\newcommand{\FR}{\mathsf{FR}}
\newcommand{\SR}{\mathsf{SR}}
\newcommand{\SecR}{\mathsf{SecR}}
\newcommand{\Sat}{\mathsf{Sat}}

\newenvironment{statusbox}[1]{%
  \def\FrameCommand{\fcolorbox{ddtablue}{ddtagray}}%
  \MakeFramed{\advance\hsize-2\width\FrameRestore}
  \noindent\textbf{#1}\par\smallskip
}{\endMakeFramed}

\newenvironment{resultbox}[1]{%
  \def\FrameCommand{\fcolorbox{ddtagreen}{ddtagray}}%
  \MakeFramed{\advance\hsize-2\width\FrameRestore}
  \noindent\textbf{#1}\par\smallskip
}{\endMakeFramed}

\newenvironment{warningbox}[1]{%
  \def\FrameCommand{\fcolorbox{ddtaorange}{ddtagray}}%
  \MakeFramed{\advance\hsize-2\width\FrameRestore}
  \noindent\textbf{#1}\par\smallskip
}{\endMakeFramed}

\newcolumntype{Y}{>{\RaggedRight\arraybackslash}X}

\begin{document}

\begin{titlepage}
\centering
\vspace*{1.6cm}
{\Large Documentation-Driven Threat Analysis\par}
\vspace{0.8cm}
{\Huge\bfseries SecurityRequirement\par}
\vspace{0.30cm}
{\LARGE S2 - chiusura semantica del metamodel e corpus facial-access camera\par}
\vspace{0.9cm}
{\large Dal SpecializedRequirement astratto alla specializzazione di sicurezza concreta\par}
\vfill
\begin{minipage}{0.86\textwidth}
\small
\textbf{Stato.} Chiusura semantica S2 per il baseline corrente. I probe hanno chiuso il delta minimo di \texttt{SecurityRequirement}; restano fuori Finding, Base Analysis, STRIDE, risk, control e il meccanismo strutturale di provenance.

\medskip
\textbf{Baseline repository.} \texttt{nballestriero/documentation-driven-threat-analysis}, commit \texttt{b9bbe690ba9f73e38f75ff9e7d8d36a88a16f49a}.

\medskip
\textbf{Obiettivo.} Formalizzare il minimo delta semantico di SecurityRequirement rispetto a SpecializedRequirement e conservarne la closure evidence sul corpus telecamera gia' usato in S1.
\end{minipage}
\vfill
{\large 14 agosto 2026\par}
\end{titlepage}

\tableofcontents
\newpage

\section{Scope S2 e baseline ereditato}

S2 non riapre i layer gia' chiusi. Assume il contratto documentale raggiunto fino a S1.5:

\begin{lstlisting}
GovernedDocument
    |
    +-- MacroRequirement
    +-- Decision
    `-- Requirement [abstract]
            |
            +-- FunctionalRequirement
            `-- SpecializedRequirement [abstract]
\end{lstlisting}

con il contratto comune:

\begin{lstlisting}
Requirement [abstract]
    extends GovernedDocument
    normativeClause : NormativeClause [1..*]
\end{lstlisting}

Ogni Requirement e' esso stesso una obbligazione normativa governata; tutte le sue clausole devono costituire una sola unita' normativa coerente. Una obbligazione indipendentemente evolvibile o assessable deve essere modellata come Requirement separato.

Per SpecializedRequirement restano chiusi: parent FR singolo, normative strengthening, composizione congiuntiva, removal test, no duplicazione della correttezza funzionale, subordinate positive action, autonomia operativa come boundary verso FR, realization separation e parent dependence.

\begin{statusbox}{S2 boundary}
S2 determina che cosa renda uno SpecializedRequirement specificamente un SecurityRequirement e quale sia il suo contratto minimo aggiuntivo. La struttura dell'analisi che lo motiva resta downstream e non deve essere incorporata nel requisito per convenienza.
\end{statusbox}

\section{Due assi da non confondere: tipo e ownership}

La specializzazione di tipo e la catena di ownership sono due viste diverse.

\subsection{Gerarchia di tipo}

\begin{figure}[h]
\centering
\begin{tikzpicture}[
  node distance=10mm and 16mm,
  box/.style={draw=ddtadark,rounded corners,align=center,inner sep=6pt,minimum width=43mm},
  arr/.style={-{Latex[length=2.4mm]},thick}
]
\node[box] (gd) {GovernedDocument};
\node[box, below=of gd] (req) {Requirement\\\textit{abstract}};
\node[box, below left=of req] (fr) {FunctionalRequirement};
\node[box, below right=of req] (sr) {SpecializedRequirement\\\textit{abstract}};
\node[box, below=of sr] (sec) {SecurityRequirement};
\draw[arr] (gd)--(req);
\draw[arr] (req)--(fr);
\draw[arr] (req)--(sr);
\draw[arr] (sr)--(sec);
\end{tikzpicture}
\caption{Gerarchia IS-A. SecurityRequirement e' un tipo concreto di SpecializedRequirement, non un documento figlio aggiuntivo.}
\end{figure}

\subsection{Catena documentale concreta}

\begin{figure}[h]
\centering
\begin{tikzpicture}[
  node distance=8mm,
  box/.style={draw=ddtadark,rounded corners,align=center,inner sep=6pt,minimum width=48mm},
  arr/.style={-{Latex[length=2.4mm]},thick}
]
\node[box] (p) {Project problem framing};
\node[box, below=of p] (mr) {MacroRequirement};
\node[box, below=of mr] (d) {Decision};
\node[box, below=of d] (fr) {FunctionalRequirement};
\node[box, below=of fr] (sec) {SecurityRequirement\\\small concrete SpecializedRequirement};
\draw[arr] (p)--(mr);
\draw[arr] (mr)--(d);
\draw[arr] (d)--(fr);
\draw[arr] (fr)--(sec);
\end{tikzpicture}
\caption{Ownership/decomposition dei documenti concreti nel caso security.}
\end{figure}

Non esiste quindi una istanza intermedia \texttt{SpecializedRequirement} tra FR e SecurityRequirement. Lo stesso oggetto concreto e' contemporaneamente SecurityRequirement, SpecializedRequirement, Requirement e GovernedDocument.

\section{Domanda fondamentale S2}

\begin{statusbox}{Fundamental question}
Dato uno SpecializedRequirement gia' valido rispetto a un singolo FunctionalRequirement, quale informazione e quale semantica aggiuntiva devono essere presenti affinche' esso sia specificamente un SecurityRequirement, senza incorporare la causa analitica, l'attacco, il finding o il controllo che ne motivano o realizzano l'obbligazione?
\end{statusbox}

La domanda separa quattro concetti:

\begin{enumerate}
  \item \textbf{obbligazione governata}: il SecurityRequirement stesso;
  \item \textbf{security property}: la proprieta' di sicurezza che il requisito intende preservare;
  \item \textbf{security failure mode}: lo stato/comportamento che renderebbe osservabile la perdita di quella proprieta' nel contesto del parent FR;
  \item \textbf{cause}: attacco, errore, guasto o altro evento capace di condurre al failure mode.
\end{enumerate}

S2 verifica quali di questi elementi appartengano alla struttura intrinseca del documento e quali debbano restare nella semantica normativa o nell'analysis model.

\section{Sanity check terminologico esterno}

La definizione DDTA non viene derivata da una tassonomia esterna, ma due osservazioni NIST sono utili come controllo di non-contraddizione:

\begin{itemize}
  \item NIST descrive il rischio di sicurezza dei sistemi in termini di perdita di confidentiality, integrity o availability e riconosce che ulteriori proprieta' come authenticity e accountability possono essere rilevanti;
  \item la metodologia di risk assessment NIST SP 800-30 contempla sorgenti/eventi adversarial e non-adversarial e dispone template distinti per entrambe le famiglie.
\end{itemize}

Queste osservazioni non impongono il nostro metamodel, ma supportano una scelta importante: non definire \emph{security} come sinonimo di \emph{attacco intenzionale}.

\begin{warningbox}{Scope discipline}
La semantica SecurityRequirement puo' essere cause-neutral senza obbligare la tesi a offrire coverage esaustivo di fault, safety, reliability o resilience. Una tecnica di analisi specifica puo' deliberatamente coprire solo una parte delle cause; la limitazione di coverage appartiene all'analysis method, non alla definizione del requisito.
\end{warningbox}

\section{Definizione semantica chiusa}

\begin{resultbox}{\closed{} - S2 definition}
Un \emph{SecurityRequirement} e' uno SpecializedRequirement la cui obbligazione normativa preserva una singola proprieta' di sicurezza governata rilevante per il parent FunctionalRequirement e rende identificabile, nelle proprie normative clause, il security failure mode che costituirebbe la perdita o violazione di tale proprieta' nel contesto del parent FR.
\end{resultbox}

La definizione mantiene quindi:

\[
  \SecR \subseteq \SR
\]

ed eredita da S1:

\[
  \mathsf{Effective}(F)=F \land S_1 \land \cdots \land S_n.
\]

Il fatto che $S_i$ sia un SecurityRequirement non modifica la semantica congiuntiva: la specializzazione Security determina il concern e il contratto aggiuntivo, non una nuova regola di composizione.

\section{Forma minima: quattro alternative}

Sono state confrontate quattro forme candidate.

\begin{longtable}{p{0.12\textwidth}p{0.31\textwidth}p{0.47\textwidth}}
\toprule
\textbf{Forma} & \textbf{Struttura} & \textbf{Problema da testare} \\
\midrule
\endhead
A & nessun nuovo campo & Property e failure mode devono essere ricostruiti interamente dalle clause. \\
\addlinespace
B & \texttt{protectedSecurityProperty [1]} & La property diventa interrogabile; failure mode resta single-source nelle clause. \\
\addlinespace
C & \texttt{securityFailureMode [1]} & Il failure mode e' strutturato, ma la classificazione/coverage per property resta implicita. \\
\addlinespace
D & property + failure mode & Massima esplicitezza, ma possibile duplicazione tra failure mode strutturato e contenuto normativo. \\
\bottomrule
\end{longtable}

I probe successivi eliminano le forme ridondanti o semanticamente insufficienti. La closure S2 seleziona la forma B: \texttt{protectedSecurityProperty [1]} e' strutturata, mentre il failure mode rimane obbligatoriamente esplicito nelle normative clause. Le forme A e C sono insufficienti per il contratto desiderato; la forma D rimane non necessaria finche' una identita' strutturale separata del failure mode non sia forzata da un controesempio.

\section{SecurityProperty come delta strutturale}

\subsection{Perche' non basta inferire sempre la property dal testo}

In un singolo documento una property come Confidentiality e' spesso ovvia leggendo la clause. DDTA deve pero' poter sostenere anche query governate quali:

\begin{lstlisting}
show all SecurityRequirements protecting Confidentiality
which SecurityProperties are represented under FR-3.4?
which requirements independently protect Integrity?
\end{lstlisting}

Se la property non ha una rappresentazione strutturata, il tool dovrebbe reinterpretare continuamente il linguaggio naturale per rispondere. Questo indebolirebbe stable identity, single-source governed knowledge e tool non-authority.

\begin{resultbox}{\closed{} - structural delta}
\begin{lstlisting}
SecurityRequirement
    extends SpecializedRequirement

    protectedSecurityProperty : SecurityProperty [1]
\end{lstlisting}
\end{resultbox}

\texttt{SecurityProperty} e' per ora una governed reference concettuale. S2 non chiude un enum universale ne' una tassonomia esaustiva.

\subsection{Perche' cardinalita' uno}

S1.5 ha chiuso il coherent-unit invariant. Se Confidentiality e Integrity possono essere introdotte, soddisfatte, revisionate o ritirate indipendentemente, non devono essere aggregate nello stesso Requirement soltanto perche' entrambe sono security-related.

\begin{lstlisting}
SR-3.4-C : SecurityRequirement
    protectedSecurityProperty = Confidentiality

SR-3.4-I : SecurityRequirement
    protectedSecurityProperty = Integrity
\end{lstlisting}

La cardinalita' uno non implica che la futura tassonomia SecurityProperty debba essere piatta o mutuamente esclusiva. Una property puo' essere raffinata o specializzata; S2 chiude soltanto un riferimento governante singolo per ogni Requirement coerente.

\section{Security failure mode: semantica obbligatoria, campo non ancora necessario}

Un SecurityRequirement povero come:

\begin{lstlisting}
Ensure confidentiality.
\end{lstlisting}

non e' sufficiente per DDTA. La property identifica il concern, ma non dice che cosa significhi concretamente perderlo nel comportamento del parent FR.

S2 richiede quindi failure-mode explicitness:

\begin{resultbox}{\closed{} - SEC-01 Failure-mode explicitness}
Le normative clause di un SecurityRequirement MUST rendere identificabile lo stato, comportamento o risultato che costituirebbe la violazione della protectedSecurityProperty nel contesto del parent FunctionalRequirement.
\end{resultbox}

Per ora non viene introdotto:

\begin{lstlisting}
securityFailureMode : SecurityFailureMode [1]
\end{lstlisting}

come campo obbligatorio separato. Nei probe correnti il failure mode coincide con informazione normativa gia' necessaria alla clause; duplicarlo come secondo source of truth non aggiunge responsabilita' semantica indipendente.

\begin{statusbox}{OPEN - structural extraction}
Un futuro analysis/tool layer potra' richiedere una rappresentazione strutturata del failure mode. S2 la introdurra' solo se un controesempio dimostra che la sola explicitness nelle clause non consente governance, analisi o tracciabilita' senza perdita.
\end{statusbox}

\section{Cause neutrality}

\subsection{Failure mode non e' causa}

La catena concettuale e':

\begin{center}
\begin{tikzpicture}[
  node distance=8mm,
  box/.style={draw=ddtadark,rounded corners,align=center,inner sep=6pt,minimum width=49mm},
  arr/.style={-{Latex[length=2.4mm]},thick}
]
\node[box] (cause) {Cause / threat event\\\small attack, error, fault, other};
\node[box, below=of cause] (fm) {Security failure mode};
\node[box, below=of fm] (prop) {Loss / violation\\of SecurityProperty};
\node[box, below=of prop] (impact) {Possible adverse consequence};
\draw[arr] (cause)--(fm);
\draw[arr] (fm)--(prop);
\draw[arr] (prop)--(impact);
\end{tikzpicture}
\end{center}

Il SecurityRequirement governa la proprieta' e il failure mode; la causa spiega come si possa arrivare a quel failure mode e appartiene al futuro analysis model.

\begin{resultbox}{\closed{} - SEC-02 Cause neutrality}
L'identita' e la validita' normativa di un SecurityRequirement MUST NOT dipendere da una singola causa adversarial o non-adversarial. Attacco, errore umano, misconfiguration, software fault o hardware fault possono motivare review dello stesso failure mode senza diventare parte necessaria dell'identita' del requisito.
\end{resultbox}

Cause neutrality non significa che tutte le cause siano equivalenti nell'analisi, nella likelihood o nei controlli. Significa soltanto che il requisito non viene ridefinito ogni volta che cambia la spiegazione causale della stessa perdita di proprieta'.

\section{Corpus facial-access camera: baseline funzionale}

Il corpus riusa la decomposizione S1 senza introdurre una nuova architettura.

\subsection{MacroRequirement}

\begin{description}
  \item[\texttt{MR-0003} - Riconoscimento facciale per la verifica della persona.] Il progetto possiede la responsabilita' macro di produrre evidenza di riconoscimento facciale usabile nel processo di verifica della persona.
\end{description}

\subsection{Decision di placement}

\begin{description}
  \item[\texttt{D-3.3} - Recognition placement.] \texttt{CameraSubsystem} acquisisce l'evidenza al punto di accesso; un \texttt{RecognitionProcessor} separato esegue localmente/remotamente rispetto alla camera il riconoscimento governato dal corpus.
\end{description}

Nel probe corrente la capture deve quindi essere consegnata dal CameraSubsystem al RecognitionProcessor.

\subsection{FunctionalRequirement}

\begin{description}
  \item[\texttt{FR-3.4} - Deliver RecognitionCapture.] Per ogni \texttt{RecognitionCapture} destinata alla valutazione remota, il \texttt{CameraSubsystem} MUST consegnare la capture al \texttt{RecognitionProcessor} associandola alla corretta \texttt{RecognitionRequest}; una consegna non completata MUST NOT essere rappresentata come completata con successo.
\end{description}

Questa e' correttezza funzionale ordinaria. Associare la capture alla richiesta sbagliata non viene spostato in SecurityRequirement soltanto perche' potrebbe essere sfruttato da un attaccante.

\section{SR-3.4-C come SecurityRequirement: Confidentiality}

\begin{lstlisting}
id:
    SR-3.4-C

type:
    SecurityRequirement

parentFunctionalRequirement:
    FR-3.4

protectedSecurityProperty:
    Confidentiality

normativeClause:
    Durante la consegna governata da FR-3.4, il contenuto
    della RecognitionCapture MUST NOT diventare intelligibile
    a un soggetto non autorizzato ad accedere a quella capture.
\end{lstlisting}

Il failure mode e' identificabile senza un campo duplicato:

\begin{lstlisting}
RecognitionCapture becomes intelligible
or disclosed to an unauthorized subject.
\end{lstlisting}

\subsection{Regression checks}

\begin{itemize}
  \item \textbf{Parent ownership}: esattamente FR-3.4 - PASS.
  \item \textbf{Removal}: senza SR-3.4-C la delivery resta funzione riconoscibile, ma perde confidentiality - PASS.
  \item \textbf{Security specificity}: property esplicita e failure mode contestualizzato - PASS.
  \item \textbf{Realization separation}: TLS, payload encryption o altra strategia equivalente non cambiano il Requirement - PASS.
\end{itemize}

\section{SR-3.4-I come SecurityRequirement: Integrity}

\begin{lstlisting}
id:
    SR-3.4-I

type:
    SecurityRequirement

parentFunctionalRequirement:
    FR-3.4

protectedSecurityProperty:
    Integrity

normativeClause:
    Una RecognitionCapture alterata dopo l'acquisizione
    MUST NOT essere accettata come la capture governata
    associata alla relativa RecognitionRequest.
\end{lstlisting}

Il failure mode e':

\begin{lstlisting}
altered RecognitionCapture is accepted
as the governed capture for the request.
\end{lstlisting}

Il requisito non prescrive hash, MAC, firma, protocollo o layer di protezione. Queste sono possibili realization choices downstream.

\section{SR-3.4-P come SecurityRequirement: Authorized Provenance}

\begin{lstlisting}
id:
    SR-3.4-P

type:
    SecurityRequirement

parentFunctionalRequirement:
    FR-3.4

protectedSecurityProperty:
    Authenticity / Authorized Provenance   [candidate label]

normativeClause 1:
    Una RecognitionCapture MUST NOT essere accettata come
    evidenza valida se non e' stabilito che la sua origine
    e' autorizzata a fornire evidenza per la relativa request.

normativeClause 2:
    Il mancato accertamento di una origine autorizzata
    MUST NOT essere rappresentato come accettazione riuscita.
\end{lstlisting}

Le due clause restano un Requirement soltanto se acceptance condition e failure outcome sono parti inseparabili della stessa obbligazione di authorized provenance.

\begin{warningbox}{OPEN - SecurityProperty vocabulary}
Il label \texttt{Authenticity / Authorized Provenance} e' deliberatamente candidato. S2 non decide ancora se Authorized Provenance sia una SecurityProperty autonoma, un refinement di Authenticity o una governed property definita dal progetto. La tassonomia non deve essere chiusa solo per accomodare il corpus.
\end{warningbox}

\section{Split test: una property per Requirement}

Consideriamo la proposta errata:

\begin{lstlisting}
SR-3.4-X
    protectedSecurityProperty = {Confidentiality, Integrity}
    clause 1 = capture MUST NOT be disclosed
    clause 2 = capture MUST NOT be altered and accepted
\end{lstlisting}

Le due obbligazioni possono essere violate, implementate, assessate e revisionate indipendentemente. Il coherent-unit invariant S1.5 forza quindi lo split:

\begin{lstlisting}
SR-3.4-C -> Confidentiality
SR-3.4-I -> Integrity
\end{lstlisting}

\begin{resultbox}{\closed{} - SEC-03 Single governing security property}
Ogni SecurityRequirement possiede esattamente un riferimento governante a SecurityProperty. Se due property esprimono obbligazioni indipendentemente evolvibili/assessable, devono essere modellate come SecurityRequirement distinti.
\end{resultbox}

\section{Cause mutation sullo stesso SecurityRequirement}

Prendiamo SR-3.4-C. Manteniamo invariati property, failure mode e normative clause, cambiando soltanto la causa ipotizzata.

\begin{longtable}{p{0.22\textwidth}p{0.30\textwidth}p{0.35\textwidth}}
\toprule
\textbf{Mutation} & \textbf{Cause} & \textbf{Esito sul requisito} \\
\midrule
\endhead
A & malicious network observation/interception & SR-3.4-C invariato; la causa e' adversarial. \\
\addlinespace
B & endpoint misconfiguration exposing content & SR-3.4-C invariato; la causa e' non-adversarial/accidentale. \\
\addlinespace
C & software defect exposing capture & SR-3.4-C invariato; cambia l'origine causale, non il failure mode. \\
\addlinespace
D & hardware replacement/disposal path exposes stored copy & SR-3.4-C puo' continuare a esprimere confidentiality se il trattamento ricade nel parent/applicability governato; eventuale cambio di parent richiede review, non auto-migrazione. \\
\bottomrule
\end{longtable}

La mutazione supporta SEC-02: l'attacco non e' prerequisito ontologico del SecurityRequirement.

\section{Caso limite: perdita dati per guasto}

La frase ``perdita dati dovuta a rottura'' non basta da sola a classificare un Requirement come security.

\subsection{Guasto con solo failure funzionale ordinario}

Se un componente si rompe e FR-3.4 non completa la delivery, la clausola funzionale gia' richiede che la consegna fallita non venga rappresentata come successo. Questo rimane FR correctness, anche se il guasto e' grave.

\subsection{Guasto che viola una property di sicurezza governata}

Se invece il progetto governa una property di Availability/Integrity e il failure mode e', per esempio, perdita permanente di un record necessario oppure indisponibilita' non tollerabile di una capability autorizzata, un evento non-adversarial puo' motivare un SecurityRequirement relativo a quella property.

Il criterio non e':

\begin{lstlisting}
Was there an attacker?
\end{lstlisting}

ma:

\begin{lstlisting}
Is there a governed SecurityProperty whose loss is made
explicit by this Requirement and whose failure mode is
separate from ordinary functional correctness?
\end{lstlisting}

Lo stesso fenomeno puo' essere simultaneamente rilevante per reliability/resilience. S2 non pretende di rendere SecurityRequirement l'unica classificazione possibile di ogni guasto.

\section{Availability probe e confine con reliability}

Per stressare il modello senza estendere il corpus canonico, consideriamo un probe illustrativo sotto un FR che processa richieste valide:

\begin{lstlisting}
protectedSecurityProperty:
    Availability

normativeClause:
    Repeated invalid submissions MUST NOT permanently starve
    governed valid RecognitionRequests from receiving the
    processing service owned by the parent FR.
\end{lstlisting}

Il failure mode e' starvation permanente delle richieste valide. Una causa potrebbe essere adversarial (denial-of-service intenzionale) oppure non-adversarial (client difettoso che genera richieste invalide). Il Requirement puo' restare semanticamente uguale.

Negative control: un generico ``the service must never fail'' non e' automaticamente SecurityRequirement; puo' essere reliability/performance e manca di una property/failure mode sufficientemente governati.

\section{Negative controls S2}

\subsection{NC-SEC-01 - Functional correctness}

\begin{lstlisting}
RecognitionCapture MUST be associated with the correct
RecognitionRequest.
\end{lstlisting}

E' necessaria alla funzione di delivery stessa. Rimane in FR-3.4 anche se una associazione errata puo' produrre conseguenze di sicurezza.

\subsection{NC-SEC-02 - Control/realization leakage}

\begin{lstlisting}
RecognitionCapture MUST be protected with TLS 1.3.
\end{lstlisting}

Nella forma corrente e' una scelta di realizzazione/Decision candidata, non la semantica minima del SecurityRequirement. Il requisito deve esprimere property e failure mode senza prescrivere incidentalmente il meccanismo.

\subsection{NC-SEC-03 - Attack description}

\begin{lstlisting}
An attacker MUST NOT perform a MITM attack.
\end{lstlisting}

Non definisce in modo sufficiente quale property del parent FR debba essere preservata, quale comportamento costituisca failure e quale obbligo governato debba valere. E' una descrizione causale/threat-oriented, non un SecurityRequirement completo.

\subsection{NC-SEC-04 - Generic quality concern}

\begin{lstlisting}
Delivery should be reliable and fast.
\end{lstlisting}

Non e' una obligation sufficientemente assessable e non identifica una SecurityProperty. Potrebbe generare futuri Performance/QualityRequirement, ma non e' SecurityRequirement per sola rilevanza operativa.

\section{Architecture mutation regression}

S2 deve conservare le mutation gia' superate da S1.

\begin{longtable}{p{0.23\textwidth}p{0.30\textwidth}p{0.35\textwidth}}
\toprule
\textbf{Mutation} & \textbf{FR/SecR} & \textbf{Interpretazione} \\
\midrule
\endhead
Ethernet $\rightarrow$ Wi-Fi & FR-3.4 e SecR possono restare invariati dopo review & Cambia il mezzo, non necessariamente l'obbligo. \\
\addlinespace
Transport external $\rightarrow$ project-owned & puo' emergere un nuovo FR di transport capability; SecR su FR-3.4 possono restare & Ownership della realizzazione non trasferisce ownership normativa. \\
\addlinespace
Recognition remote $\rightarrow$ local to camera & FR-3.4 scompare; i SecR figli vengono riesaminati e possono diventare non-applicable/retired & Conferma parent dependence e no auto-migration. \\
\addlinespace
Raw capture $\rightarrow$ opaque reference & review locale dei SecR interessati & Il modello richiede effective governed context, ma non introduce in S2 una generic \texttt{affects} relation. \\
\bottomrule
\end{longtable}

\section{Whole-chain validation}

Il corpus consente ora di attraversare l'intero metamodel della documentazione governata fino alla sicurezza:

\begin{figure}[h]
\centering
\begin{tikzpicture}[
  node distance=7mm,
  box/.style={draw=ddtadark,rounded corners,align=center,inner sep=6pt,minimum width=55mm},
  sec/.style={draw=ddtablue,rounded corners,align=center,inner sep=5pt,minimum width=40mm,font=\small},
  arr/.style={-{Latex[length=2.4mm]},thick}
]
\node[box] (p) {Project framing\\controlled facial-access};
\node[box, below=of p] (mr) {MR-0003\\Facial recognition for person verification};
\node[box, below=of mr] (d) {D-3.3\\Recognition placement};
\node[box, below=of d] (fr) {FR-3.4\\Deliver RecognitionCapture};
\node[sec, below left=8mm and 13mm of fr] (c) {SR-3.4-C\\SecurityRequirement\\Confidentiality};
\node[sec, below=8mm of fr] (i) {SR-3.4-I\\SecurityRequirement\\Integrity};
\node[sec, below right=8mm and 13mm of fr] (p2) {SR-3.4-P\\SecurityRequirement\\Authorized provenance};
\draw[arr] (p)--(mr);
\draw[arr] (mr)--(d);
\draw[arr] (d)--(fr);
\draw[arr] (fr)--(c);
\draw[arr] (fr)--(i);
\draw[arr] (fr)--(p2);
\end{tikzpicture}
\caption{Catena documentale concreta. I tre SecurityRequirement sono SpecializedRequirement concreti indipendenti dello stesso FR.}
\end{figure}

Il passaggio non richiede Threat, Attack, Finding, Control o AnalysisRecord dentro la semantica dei documenti. Questi concetti potranno spiegare coverage, motivazione, plausibilita', risk e realization downstream.

\section{Cosa e' completato con la closure S2}

Con la closure degli invarianti S2, e' dichiarato completo per il baseline corrente il \textbf{governed documentation metamodel through SecurityRequirement}:

\begin{lstlisting}
Project problem framing [method precondition]
    -> MacroRequirement
        -> Decision
            -> FunctionalRequirement
                -> SpecializedRequirement [abstract]
                    -> SecurityRequirement
\end{lstlisting}

Questo non equivale ancora al full DDTA metamodel. Restano da formalizzare almeno il lato analitico e la sua integrazione:

\begin{lstlisting}
Governed documentation
    -> Base Analysis / analyzable representation
    -> Analysis execution
    -> Finding / analytical result
    -> governed feedback / provenance
    -> requirement or other document change
\end{lstlisting}

La separazione e' intenzionale: S2 chiude la destinazione normativa senza anticipare la struttura che produrra' o motivara' i cambiamenti.

\section{Stato epistemico S2 dopo il corpus}

\begin{longtable}{p{0.20\textwidth}p{0.70\textwidth}}
\toprule
\textbf{Stato} & \textbf{Elementi} \\
\midrule
\endhead
\closed & Baseline MR/Decision/FR; Requirement abstraction S1.5; SpecializedRequirement S1; parent singolo; coherent-unit; provenance distinta dalla semantica. \\
\addlinespace
\closed & \texttt{SecurityRequirement IS-A SpecializedRequirement}; definizione property-centered; \texttt{protectedSecurityProperty : SecurityProperty [1]}; failure-mode explicitness; cause neutrality; single governing security property. \\
\addlinespace
\candidate & Forma finale del label/property per Authorized Provenance; possibilita' di usare SecurityProperty come governed reference estensibile; availability probe come generality evidence. \\
\addlinespace
\openstatus & Tassonomia/gerarchia di SecurityProperty; eventuale futura rappresentazione strutturata di SecurityFailureMode; exact L2 representation; effective governed context. \\
\addlinespace
\deferred & Attack, ThreatActor, ThreatEvent, Incident, Vulnerability, Finding, Risk, Control, Base Analysis/BAE, STRIDE formalization, analysis coverage, provenance/change-event model, verification evidence. \\
\bottomrule
\end{longtable}

\section{Invarianti S2 chiusi}

\begin{enumerate}[label=\textbf{S2-CLOSED-\arabic*.},leftmargin=*,align=left,labelsep=0.6em]
  \item \textbf{Subtype.} Ogni SecurityRequirement e' uno SpecializedRequirement e ne conserva tutti gli invarianti.
  \item \textbf{Security-property anchoring.} Ogni SecurityRequirement riferisce esattamente una SecurityProperty governante.
  \item \textbf{Failure-mode explicitness.} Le normative clause rendono identificabile che cosa costituisce perdita/violazione della property nel parent FR.
  \item \textbf{Cause neutrality.} L'identita' normativa non dipende da una specifica causa adversarial o non-adversarial.
  \item \textbf{Attack non-necessity.} La presenza di un attacco noto non e' prerequisito per la validita' del requisito.
  \item \textbf{Single coherent security obligation.} Property e clause devono formare una sola unita' normativa; property indipendenti richiedono Requirement distinti.
  \item \textbf{Functional preservation.} La rimozione del SecurityRequirement lascia riconoscibile il parent FR, pur perdendo la property di sicurezza.
  \item \textbf{No functional theft.} Una condizione di ordinaria correttezza funzionale non viene spostata in SecurityRequirement per sola security relevance.
  \item \textbf{No attack-as-requirement.} Descrivere attacker, technique o threat event non sostituisce property, failure mode e obligation.
  \item \textbf{Realization separation.} Control, protocollo, algoritmo, componente o layer non sono parte necessaria della semantica SecurityRequirement.
  \item \textbf{Method neutrality.} La semantica SecurityRequirement non richiede STRIDE o altra tecnica specifica.
  \item \textbf{Provenance separation.} La causa analitica che motiva il requisito non viene incorporata come singolo analysis id intrinseco.
\end{enumerate}

\section{Reopen criteria e confine successivo}

S2 e' chiuso per il baseline corrente e deve essere riaperto se emerge almeno uno dei seguenti controesempi:

\begin{itemize}
  \item un SecurityRequirement valido richiede necessariamente due SecurityProperty indipendenti nello stesso Requirement;
  \item un failure mode non puo' essere governato senza una identita' strutturale separata dalle normative clause;
  \item la causa adversarial/non-adversarial modifica inevitabilmente l'identita' normativa anche mantenendo invariati property e failure outcome;
  \item il modello non riesce a distinguere in modo stabile security da ordinary functional correctness o da quality/reliability;
  \item un caso security concreto forza Attack, Finding, Risk o Control dentro il core del requisito.
\end{itemize}

\begin{statusbox}{Closure boundary}
S2 chiude il governed documentation metamodel through SecurityRequirement. I successivi microstep non devono riaprire questa semantica salvo controesempio concreto. La tassonomia SecurityProperty, la possibile identita' strutturale del failure mode e il lato analitico DDTA restano rispettivamente OPEN o DEFERRED secondo lo stato epistemico sopra.
\end{statusbox}

\section{Riferimenti terminologici di controllo}

\begin{itemize}
  \item NIST SP 800-30 Rev. 1, \emph{Guide for Conducting Risk Assessments}: threat sources/events e template di assessment adversarial e non-adversarial.
  \item NIST CSRC Glossary, \emph{Information System-Related Security Risk}: rischio derivante dalla perdita di confidentiality, integrity o availability.
  \item NIST CSRC Glossary, \emph{Security}: confidentiality, integrity e availability come proprieta' centrali; ulteriori property possono essere rilevanti a seconda del contesto.
\end{itemize}

Questi riferimenti sono usati solo come sanity check terminologico. Le cardinalita', i boundary e gli invarianti DDTA del presente documento derivano dai probe del metamodel e non sono attribuiti a NIST.

\end{document}
